\section{Bước 5: Kiểm thử hệ thống bằng ping, traceroute và các công cụ khác}
\label{sec:step5}

Kết quả trong phần này lấy từ nhật ký kiểm thử của nhóm và các ảnh trong thư mục \texttt{graphics/tc\_pic}. Mục nào không có ảnh bằng chứng được đánh dấu \textbf{chưa kiểm thử}; báo cáo không ghi PASS cho ca chưa có bằng chứng. Phần bằng chứng bảo mật và VPN (SEC, DMZ, INT, VPN) được trình bày chi tiết ở mục~\ref{sec:security_test_evidence}.

\subsection{Sáu kiểm tra bắt buộc của đề bài}
\begin{small}
\setlength{\tabcolsep}{3pt}
\begin{longtable}{|c|p{2.3cm}|p{3.0cm}|p{4.6cm}|p{2.6cm}|}
\caption{Sáu kiểm tra kết nối bắt buộc} \label{tab:six_checks} \\
\hline
\textbf{\#} & \textbf{Yêu cầu} & \textbf{Lệnh} & \textbf{Kết quả ghi nhận} & \textbf{Trạng thái} \\ \hline
\endhead
1 & PC cùng VLAN & \texttt{ping} giữa hai PC cùng VLAN & Chưa có ảnh riêng & \textbf{Chưa kiểm thử} \\ \hline
2 & PC khác VLAN & \texttt{ping} từ \texttt{MAIN-PC-STAFFA} (VLAN 10) tới máy VLAN 20 và đích khác VLAN & Ảnh LAN-01: \texttt{MAIN-PC-STAFFA} ping \texttt{10.10.20.50} nhận 4/4; ảnh LAN-02: \texttt{DBP-PC1} ping \texttt{10.20.30.14} (lần đầu mất 1/4 gói, lần sau 4/4) & PASS (theo ảnh LAN-01, LAN-02) \\ \hline
3 & Main $\leftrightarrow$ chi nhánh & \texttt{ping}, \texttt{tracert 10.20.40.10} & Ảnh WAN-01 đến WAN-03: ping giữa \texttt{MAIN-PC-STAFFA}, \texttt{DBP-PC1}, \texttt{BHTQ-PC3} và máy ở site khác đều nhận 4/4; \texttt{tracert} 4 hop qua \texttt{BHTQ-R1} $\to$ \texttt{Carrier-R} $\to$ \texttt{DBP-R1} $\to$ server DBP; TC-06: 5/5, RTT 2\,ms & PASS \\ \hline
4 & Tới server DMZ & Mở \texttt{http://10.10.50.11} từ nhân viên & DMZ-01: trang hiển thị; ảnh cho thấy \texttt{10.10.50.11} là \texttt{DMZ-SERVER}, khác \texttt{DMZ-WEB} (\texttt{10.10.50.10}) & PASS \\ \hline
5 & \textbf{Khách không vào được LAN} & Ping từ \texttt{MAIN-B-GUEST} tới \texttt{10.10.20.10}, \texttt{10.10.40.10}, \texttt{10.10.50.10} & Chặn được staff và Server Farm, \textbf{nhưng ping được DMZ \texttt{10.10.50.10}} & \textbf{FAIL} (SEC-01) \\ \hline
6 & Internet tới Web server & Mở \texttt{http://192.0.2.10} từ nhân viên & INT-01: trang hiển thị & PASS \\ \hline
\end{longtable}
\end{small}

\subsection{Kiểm thử bảo mật và VPN}
Bảng~\ref{tab:security_test_summary} (mục~\ref{sec:security_test_evidence}) liệt kê các ca SEC-01, SEC-02, DMZ-01, INT-01, VPN-S2S-01, VPN-RA-01 và VPN-RA-02 cùng ảnh minh chứng. Tóm tắt:
\begin{itemize}[noitemsep,topsep=2pt]
    \item \textbf{SEC-01 (cô lập Guest): FAIL.} \texttt{MAIN-B-GUEST} bị chặn khi ping \texttt{10.10.20.10} và \texttt{10.10.40.10} nhưng vẫn ping được DMZ \texttt{10.10.50.10} (4/4) và máy ngoài \texttt{192.0.2.10} (3/4).
    \item \textbf{SEC-02:} từ \texttt{REMOTE} ping \texttt{10.10.40.10} (không qua VPN) bị timeout 4/4; PASS có giới hạn.
    \item \textbf{VPN:} IPsec site-to-site kết thúc tại \texttt{MAIN-WAN-R2} (peer \texttt{10.255.0.18}, \texttt{10.255.0.22}); VPN truy cập từ xa kết thúc tại \texttt{MAIN-VPN}, client \texttt{REMOTE} nhận \texttt{10.10.200.10}.
    \item \textbf{Chưa kiểm thử:} SSH từ Jump Host \texttt{MAIN-SSH}; Internet $\to$ DMZ qua NAT \texttt{203.162.10.3}; camera/NVR; IDS/IPS.
\end{itemize}
\noindent \textbf{Ghi chú SEC-01:} \texttt{ACL-GUEST-ISOLATE} đã được sửa trong báo cáo nhưng \textbf{chưa áp dụng vào mô phỏng và chạy lại}. Lỗi nằm ở chỗ \texttt{DMZ-WEB} \texttt{10.10.50.10} thuộc \texttt{10.0.0.0/8}; ACL phải \texttt{deny} trước khi \texttt{permit} web ra ngoài. Phải áp dụng và chạy lại trước khi ghi PASS.
\begin{lstlisting}[language=csh,caption={ACL-GUEST-ISOLATE đã sửa (chưa áp dụng vào mô phỏng)}]
ip access-list extended ACL-GUEST-ISOLATE
 permit udp any eq bootpc any eq bootps
 permit udp any host 10.10.40.10 eq domain
 permit tcp 10.10.60.0 0.0.3.255 host 10.10.40.10 eq domain
 deny   ip  10.10.60.0 0.0.3.255 10.0.0.0 0.255.255.255
 permit tcp 10.10.60.0 0.0.3.255 any eq www
 permit tcp 10.10.60.0 0.0.3.255 any eq 443
!
interface Vlan60
 ip access-group ACL-GUEST-ISOLATE in
\end{lstlisting}

\subsection{Kịch bản demo}

\subsubsection{Giai đoạn 1: Trạng thái chuẩn và hội tụ}

\paragraph{Bước 1.1: Kiểm tra HSRP.} Truy cập CLI của các thiết bị gateway chính và phụ để kiểm tra trạng thái nhóm HSRP. Hai ảnh HA-01 và HA-01.1 minh chứng cho hai switch lõi; listing cho \texttt{DBP-R1} chưa có ảnh riêng.
\begin{lstlisting}[language=csh,caption={Kiểm tra HSRP trên MAIN-CORE-SW1 và router chi nhánh}]
MAIN-CORE-SW1# show standby brief
                   P indicates configured to preempt.
Interface   Grp  Pri P State   Active          Standby         Virtual IP
Vl10        10   200 P Active  local           10.10.10.3      10.10.10.1
Vl20        20   200 P Active  local           10.10.20.3      10.10.20.1
Vl30        30   200 P Active  local           10.10.30.3      10.10.30.1
Vl40        40   200 P Active  local           10.10.40.3      10.10.40.1
Vl60        60   200 P Active  local           10.10.60.3      10.10.60.1
Vl70        70   200 P Active  local           10.10.70.3      10.10.70.1
Vl80        80   200 P Active  local           10.10.80.3      10.10.80.1
Vl90        90   200 P Active  local           10.10.90.3      10.10.90.1
Vl250       250  200 P Active  local           10.10.250.3     10.10.250.1

DBP-R1# show standby brief
Interface   Grp  Pri P State   Active          Standby         Virtual IP
Gi0/0.10    10   110 P Active  local           10.20.10.3      10.20.10.1
Gi0/0.30    30   110 P Active  local           10.20.30.3      10.20.30.1
Gi0/0.40    40   110 P Active  local           10.20.40.3      10.20.40.1
Gi0/0.70    70   110 P Active  local           10.20.70.3      10.20.70.1
Gi0/0.90    90   110 P Active  local           10.20.90.3      10.20.90.1
Gi0/0.99    99   110 P Active  local           10.20.99.3      10.20.99.1
\end{lstlisting}
\noindent \textbf{Tiêu chí nghiệm thu:} toàn bộ nhóm trên \texttt{MAIN-CORE-SW1}, \texttt{DBP-R1}, \texttt{BHTQ-R1} ở trạng thái \texttt{Active} và hiển thị đúng địa chỉ Standby của thiết bị phụ (\texttt{MAIN-CORE-SW2}, \texttt{DBP-R2}, \texttt{BHTQ-R2}); không còn nhóm Speak/Unknown.

\tcfig{HA_01.png}{HA-01: trạng thái HSRP trên MAIN-CORE-SW1 và MAIN-CORE-SW2}{fig:demo-ha-01}
\tcfig{HA_01.1.png}{HA-01.1: sơ đồ và ping liên tục từ MAIN-PC-STAFFA tới 10.10.50.12, có gián đoạn rồi phục hồi}{fig:demo-ha-01-1}

\paragraph{Bước 1.2: Kiểm tra láng giềng OSPF.}
\begin{lstlisting}[language=csh,caption={Láng giềng OSPF Area 0}]
Carrier-R# show ip ospf neighbor
Neighbor ID     Pri   State           Dead Time   Address         Interface
1.1.1.1           1   FULL/  -        00:00:32    10.255.0.2      Serial0/3/0
1.1.1.2           1   FULL/  -        00:00:35    10.255.0.6      Serial0/3/1
3.3.3.1           1   FULL/  -        00:00:31    10.255.0.10     Serial0/3/2
4.4.4.1           1   FULL/  -        00:00:33    10.255.0.14     Serial0/3/3

BHTQ-R2# show ip ospf neighbor
Neighbor ID     Pri   State           Dead Time   Address         Interface
4.4.4.1           1   FULL/DR         00:00:38    10.30.99.2      GigabitEthernet0/0.99
\end{lstlisting}
\noindent \textbf{Tiêu chí nghiệm thu:} liên kết Serial đạt \texttt{FULL/-}; liên kết subinterface đạt \texttt{FULL/DR} hoặc \texttt{FULL/BDR}.

\paragraph{Bước 1.3: Ping và truy vết đường chính.} Từ \texttt{BHTQ-PC1} (\texttt{10.30.10.10}):
\begin{lstlisting}[language=csh,caption={Ping và traceroute từ chi nhánh BHTQ}]
C:\> ping 10.10.40.10
Reply from 10.10.40.10: bytes=32 time=4ms TTL=124
Reply from 10.10.40.10: bytes=32 time=3ms TTL=124
Reply from 10.10.40.10: bytes=32 time=3ms TTL=124
Reply from 10.10.40.10: bytes=32 time=3ms TTL=124
Ping statistics for 10.10.40.10: Packets: Sent = 4, Received = 4, Lost = 0 (0% loss)

C:\> tracert 10.20.40.10
  1     1 ms     1 ms     1 ms   10.30.10.2      (BHTQ-R1)
  2     2 ms     2 ms     2 ms   10.255.0.13     (Carrier-R)
  3     3 ms     2 ms     3 ms   10.255.0.10     (DBP-R1)
  4     4 ms     3 ms     4 ms   10.20.40.10     (Server DBP)
Trace complete.
\end{lstlisting}
\noindent Các ảnh LAN-01, LAN-02 và WAN-01 đến WAN-03 dưới đây ghi nhận ping từ các máy trạm; kết quả chỉ áp dụng cho đích và lần thử hiển thị.

\tcfig{LAN_01.png}{LAN-01: MAIN-PC-STAFFA truy cập một số đích LAN và DMZ}{fig:demo-lan-01}
\tcfig{LAN_02.png}{LAN-02: DBP-PC1 ping 10.20.30.14, lần đầu mất gói rồi nhận đủ phản hồi}{fig:demo-lan-02}
\tcfig{WAN_01.png}{WAN-01: MAIN-PC-STAFFA ping 10.30.30.10 và 10.20.10.18, đều nhận đủ 4/4 phản hồi}{fig:demo-wan-01}
\tcfig{WAN_02.png}{WAN-02: BHTQ-PC3 ping 10.10.10.50 nhận đủ 4/4 phản hồi và cấu hình IP của máy}{fig:demo-wan-02}
\tcfig{WAN_03.png}{WAN-03: DBP-PC1 ping 10.20.30.14 (lần đầu mất 1/4 gói) và 10.10.10.50 nhận đủ phản hồi}{fig:demo-wan-03}

\subsubsection{Giai đoạn 2: An ninh phân vùng, DMZ và Internet}
\begin{itemize}[noitemsep,topsep=2pt]
    \item \textbf{Cô lập Guest:} từ \texttt{MAIN-B-GUEST} (VLAN 60, \texttt{10.10.60.50}), ping tới \texttt{10.10.20.10} và \texttt{10.10.40.10} không có phản hồi; ping \texttt{10.10.50.10} nhận đủ phản hồi ở lần thử sau; ping \texttt{192.0.2.10} nhận 3/4 phản hồi. Kết luận: ca cô lập Guest \textbf{FAIL} (ảnh \texttt{SEC\_01.png}).
    \item \textbf{Web DMZ:} từ \texttt{MAIN-PC-STAFFA}, mở \texttt{http://10.10.50.11} (\texttt{DMZ-SERVER}) hiển thị trang Cisco Packet Tracer (ảnh \texttt{DMZ\_01.png}).
    \item \textbf{Web Internet:} mở \texttt{http://192.0.2.10} (\texttt{INTERNET WEB}) hiển thị thành công (ảnh \texttt{INT\_01.png}).
\end{itemize}
Các ảnh này không kiểm chứng truy cập Web từ Internet vào địa chỉ NAT công khai hoặc SSH quản trị; các phép thử đó \textbf{chưa được thực hiện}.

\subsubsection{Giai đoạn 3: Chuyển đổi dự phòng WAN tại chi nhánh BHTQ}
\begin{enumerate}[noitemsep,topsep=2pt]
    \item \textbf{Luồng giám sát:} trên \texttt{BHTQ-PC3} chạy \texttt{ping 10.20.40.10 -t} tới server DBP.
    \item \textbf{Bơm sự cố đứt leased line:} trên \texttt{BHTQ-R1} ngắt cổng nối \texttt{Carrier-R}:
\begin{lstlisting}[language=csh]
BHTQ-R1(config)# interface Serial0/3/0
BHTQ-R1(config-if)# shutdown
%LINK-5-CHANGED: Interface Serial0/3/0, changed state to administratively down
%LINEPROTO-5-UPDOWN: Line protocol on Interface Serial0/3/0, changed state to down
\end{lstlisting}
    \item Bấm \texttt{Alt + D} để đẩy nhanh chu kỳ phát hiện lỗi trong mô phỏng.
    \item \textbf{Phản ứng của hệ thống:}
    \begin{itemize}[noitemsep,topsep=2pt]
        \item \textit{HSRP tracking:} \texttt{standby track Serial0/3/0 decrement 60} làm priority của \texttt{BHTQ-R1} giảm từ $110$ xuống $50$ trên toàn bộ nhóm VLAN.
        \item \textit{Preemption:} \texttt{BHTQ-R2} (priority $105 > 50$) chuyển sang \texttt{Active}.
        \item \textit{Floating static:} OSPF giữa \texttt{BHTQ-R1} và \texttt{Carrier-R} mất, tuyến \texttt{ip route 0.0.0.0 0.0.0.0 10.255.0.21 150} vào bảng định tuyến.
    \end{itemize}
    \item \textbf{Kết quả ghi nhận:} chỉ rơi \textbf{1 gói}, RTT tăng từ $3\,\text{ms}$ lên $11\,\text{ms}$ do đi qua đường broadband; \texttt{tracert} đi qua \texttt{BHTQ-R2} $\to$ \texttt{ISP2-R} $\to$ \texttt{DBP-R2} $\to$ server DBP.
\begin{lstlisting}[language=csh]
Reply from 10.20.40.10: bytes=32 time=3ms TTL=124
Request timed out.
Reply from 10.20.40.10: bytes=32 time=12ms TTL=123
Reply from 10.20.40.10: bytes=32 time=11ms TTL=123

C:\> tracert 10.20.40.10
  1     1 ms     1 ms     1 ms   10.30.10.3      (BHTQ-R2 - Active gateway moi)
  2     5 ms     4 ms     5 ms   10.255.0.21     (ISP2-R)
  3     8 ms     7 ms     7 ms   10.255.0.18     (DBP-R2)
  4     9 ms     8 ms     9 ms   10.20.40.10     (Server DBP)
\end{lstlisting}
    \item \textbf{Failback:} \texttt{no shutdown} trên \texttt{Serial0/3/0}; \texttt{BHTQ-R1} lấy lại \texttt{Active} nhờ \texttt{standby preempt}, OSPF \texttt{FULL} trở lại với \texttt{Carrier-R}.
\end{enumerate}
Các ảnh WAN-01 đến WAN-03, HA-02.1 và HA-02.2 chỉ hiển thị ping và CLI; riêng HA-02.1/HA-02.2 là CLI của \texttt{DBP-R2}. Các ảnh này \textbf{không tự chứng minh} đường dự phòng cụ thể hoặc thời gian gián đoạn.

\tcfig{HA_02.1.png}{HA-02.1: log HSRP/OSPF khi khởi động và ping 10.10.40.10 từ DBP-R2 (lần đầu 3/5, lần sau 5/5)}{fig:demo-ha-02-1}
\tcfig{HA_02.2.png}{HA-02.2: log chuyển trạng thái HSRP và ping từ DBP-R2 tới 10.10.40.10}{fig:demo-ha-02-2}

\subsubsection{Giai đoạn 4: Dự phòng lõi tại Main Site}
\begin{small}
\setlength{\tabcolsep}{4pt}
\begin{longtable}{|c|p{5.0cm}|p{4.6cm}|p{2.4cm}|}
\caption{Kiểm thử dự phòng lõi} \label{tab:core_ha} \\
\hline
\textbf{Thử nghiệm} & \textbf{Thao tác} & \textbf{Kết quả theo báo cáo} & \textbf{Bằng chứng} \\ \hline
\endhead
4.1 & \texttt{shutdown} \texttt{Fa0/2} trên \texttt{MAIN-CORE-SW1} (một cáp trong Port-channel 1), ping \texttt{10.10.40.10 -t} từ \texttt{MAIN-PC-STAFFA} & 0\% mất gói, LACP dồn lưu lượng sang cáp còn lại & Mô tả trong báo cáo, chưa có ảnh \\ \hline
4.2 & Tắt \texttt{MAIN-CORE-SW1} & \texttt{MAIN-CORE-SW2} thành \texttt{Active} (holdtime 10\,s); đường VLAN 253 lên \texttt{MAIN-FW-01} Gi1/5 được thiết kế làm tuyến dự phòng & \textbf{Chưa kiểm chứng}: liên kết down trên sơ đồ, chưa có ảnh \\ \hline
\end{longtable}
\end{small}
\noindent Khi chạy thử 4.2 cần bỏ \texttt{shutdown} trên \texttt{Gi1/5} của \texttt{MAIN-FW-01} (cấu hình hiện đang shutdown). Kết quả mong đợi của 4.2 là hoạt động truy cập dịch vụ Server Farm vẫn liên tục, nhưng điều này chưa được xác nhận.

\subsection{Bảng nghiệm thu TC-01 đến TC-06}
\begin{small}
\setlength{\tabcolsep}{3pt}
\begin{longtable}{|c|p{2.5cm}|p{2.2cm}|c|p{5.4cm}|}
\caption{Bảng kết quả kiểm thử TC-01 đến TC-06} \label{tab:test_matrix} \\
\hline
\textbf{Mã} & \textbf{Hạng mục} & \textbf{Ban đầu} & \textbf{Kết quả} & \textbf{Nguyên nhân và xử lý} \\ \hline
\endfirsthead
\hline \textbf{Mã} & \textbf{Hạng mục} & \textbf{Ban đầu} & \textbf{Kết quả} & \textbf{Nguyên nhân và xử lý} \\ \hline
\endhead
TC-01 & Định tuyến \texttt{BHTQ-R2} tới server DBP \texttt{10.20.40.10} & Mất 100\% (0/5) & PASS & Default static AD 1 đè OSPF gây bất đối xứng. Đổi thành floating AD 150, bỏ \texttt{default-information originate}, đặt \texttt{passive-interface} trên cổng broadband. \\ \hline
TC-02 & Láng giềng OSPF qua VLAN 99 (BHTQ) & Kẹt INIT/DROTHER & PASS & Lệch thẻ do \texttt{dot1Q native} kết hợp HSRP/OSPF. Bỏ \texttt{native}, dùng \texttt{dot1Q 99}, native VLAN switch về 1, xóa tiến trình OSPF. \\ \hline
TC-03 & Bảng MAC trên trunk VLAN 99 & Học đủ MAC của R1 và R2 & PASS & L2 ổn định, không vòng STP, không khóa Port-Security. \\ \hline
TC-04 & Đồng bộ HSRP nhóm 90 và 99 & Nhóm 90 báo Speak & PASS & Thiếu \texttt{Gi0/0.90}, \texttt{Gi0/0.99} và allowed VLAN trên trunk. Khai báo đồng bộ trên R1/R2, mở trunk VLAN 1, 10, 30, 40, 70, 90, 99. \\ \hline
TC-05 & Traceroute DBP sang server BHTQ \texttt{10.30.40.10} & Rớt hop 4 & PASS & Server thiếu gateway. Đặt \texttt{10.30.40.1}, gán access VLAN 40. \\ \hline
TC-06 & Ping trực tiếp qua leased line & 5/5, RTT 2\,ms & PASS & Tuyến chính qua \texttt{Carrier-R} ổn định, được ưu tiên khi không có sự cố. \\ \hline
\end{longtable}
\end{small}
\noindent Chỉ các cấu hình đã xuất (15 thiết bị) mới được xác nhận là đã lưu; chưa chứng minh đã sao lưu đủ toàn bộ thiết bị hoặc đã kiểm thử đầy đủ ACL/NAT. Bằng chứng khắc phục native VLAN hiện ghi nhận tại BHTQ; cấu hình tương ứng tại DBP cần kiểm tra riêng.

\subsection{Việc cần làm để hoàn thiện kiểm thử}
\begin{enumerate}[noitemsep,topsep=2pt]
    \item Sửa và chạy lại SEC-01 (áp dụng ACL vào \texttt{Vlan60}).
    \item Chạy thêm: ping cùng VLAN, Internet $\to$ DMZ qua \texttt{203.162.10.3}, SSH từ \texttt{MAIN-SSH}, camera $\to$ NVR, thử nghiệm 4.2.
    \item Chụp ảnh từng ca với mã trùng bảng trong báo cáo (ví dụ \texttt{LAN-01}).
    \item Không ghi PASS cho ca chưa có ảnh bằng chứng.
\end{enumerate}
